\subsection{Conclusion}

Perching drones have potential to improve the efficiency of infrastructure inspection, make remote and dangerous 
ecosystems more accessible for monitoring and study, and enable safer and faster data collection and communication
in disaster responses. Recent research has produced many different perching concepts, though they still seem to be 
difficult for researchers to put into practice. The design proposed in this thesis is simple to fabricate and adapt
for different drone designs and use cases, and proven to support drones exceeding 1kg in mass.
The detection method proposed can be implemented with affordable 
components and free, open-sourced software that many developers will be familiar with. It achieves accuracy of under 75mm error in distance
estimation and reliability of over 80\% in favorable conditions, though work remains to improve robustness of detection against 
complex background patterns and stability of the detected perch orientation.
Though the prototype designed here is limited to idealized environments, it demonstrates that the proposed design is suitable for
identifying and attaching to perches, and has promise for future development. 